\documentclass[12pt, a4paper]{article}
\usepackage[utf8]{inputenc}
\usepackage{geometry}
\geometry{a4paper, margin=1in}
\usepackage{graphicx}
\usepackage{setspace}
\usepackage{titlesec}
\usepackage{booktabs}
\usepackage{amsmath}
\usepackage{amssymb}
\usepackage{enumitem}
\usepackage{pgfgantt}
\usepackage[hidelinks]{hyperref}

\setstretch{1.2}

\titleformat{\section}{\large\bfseries}{\thesection}{1em}{}
\titleformat{\subsection}{\normalsize\bfseries}{\thesubsection}{1em}{}
\titleformat{\subsubsection}{\normalsize\itshape}{\thesubsubsection}{1em}{}

\begin{document}

% ----------------------------------------------------------------------
% TITLE PAGE
% ----------------------------------------------------------------------
\begin{titlepage}
    \centering
    
    {\Large \textbf{PROJECT SYNOPSIS}} \\[0.3cm]
    {\large \textbf{B.Tech Computer Science \& Engineering --- Major Project}} \\[0.5cm]
    {\normalsize \textbf{Session 2026--27}} \\[1cm]
    
    \rule{\linewidth}{0.5mm} \\[0.3cm]
    {\LARGE \textbf{Task-Conditioned Evaluation of Quantized Large Language Models:}} \\[0.2cm]
    {\Large \textbf{Accuracy Degradation, the Reliability of Perplexity as a Proxy, and Practical Bit-Budget Guidelines}} \\[0.3cm]
    \rule{\linewidth}{0.5mm} \\[0.8cm]
    
    \textbf{Group ID:} DTU/CSE/BTP/GP159 \\[0.5cm]
    
    \textbf{Submitted by:} \\[0.2cm]
    \begin{tabular}{lll}
        \textbf{S. No.} & \textbf{Name} & \textbf{Enrollment No.} \\
        1 & Ekansh Agarwal & 23/CS/149 \\
        2 & Faqre Alam & 23/CS/150 \\
        3 & Guneet Toppo & 23/CS/160
    \end{tabular} \\[0.8cm]
    
    \textbf{Under the Supervision of:} \\[0.2cm]
    \textbf{Ms. Gull Kaur} \\ 
    Assistant Professor, Department of Computer Science \& Engineering
    
    \vfill
    
    \includegraphics[width=3.5cm]{dtu_logo_hd.png} \\[0.3cm]
    
    {\large \textbf{Department of Computer Science \& Engineering}} \\[0.1cm]
    {\Large \textbf{Delhi Technological University}} \\[0.1cm]
    {\normalsize (Formerly Delhi College of Engineering)} \\[0.1cm]
    {\normalsize Bawana Road, Delhi - 110042}
    
\end{titlepage}

\newpage
\tableofcontents
\newpage

% ----------------------------------------------------------------------
% 1. PROJECT TITLE & GROUP DETAILS
% ----------------------------------------------------------------------
\section{Project Title and Group Details}

\subsection{Project Title}
\textbf{Task-Conditioned Evaluation of Quantized Large Language Models: Accuracy Degradation, the Reliability of Perplexity as a Proxy, and Practical Bit-Budget Guidelines}

\subsection{Group Details}
\begin{itemize}[leftmargin=*]
    \item \textbf{Group ID:} DTU/CSE/BTP/GP159
    \item \textbf{Guide:} Ms. Gull Kaur, Assistant Professor, Department of Computer Science \& Engineering, DTU
    \item \textbf{Group Members:}
    \begin{enumerate}
        \item Ekansh Agarwal (23/CS/149)
        \item Faqre Alam (23/CS/150)
        \item Guneet Toppo (23/CS/160)
    \end{enumerate}
    \item \textbf{Junior Team Members (2nd/3rd Year):} [To be filled]
\end{itemize}

% ----------------------------------------------------------------------
% 2. ABSTRACT AND PROBLEM STATEMENT
% ----------------------------------------------------------------------
\section{Abstract and Problem Statement}

Large language models (LLMs) are usually stored and run in quantized form, which means the weights are compressed from 16 bits down to 8 or 4 bits so the model requires significantly less memory. Early foundational papers like GPTQ and AWQ showed that 4-bit models lose very little accuracy. Jin et al. later tested quantized models across ten benchmarks and concluded that perplexity, which is computationally cheap to calculate, is a decent indicator of how well a model will perform on downstream tasks.

There are two major problems with this prevailing paradigm:
\begin{enumerate}[leftmargin=*]
    \item Most of that empirical evidence was collected before multi-step reasoning benchmarks such as GSM8K and instruction-following benchmarks like IFEval became standard. A model that must execute several chained logical steps may react to quantization completely differently from one that only has to recall a static fact. Recent 2026 studies suggest the picture is significantly more complicated for reasoning tasks; however, they focus exclusively on models of 3B parameters or more and examine only a single quantization format.
    \item Small models in the 0.5B to 1B parameter range---which represent the practical models that students, edge developers, and consumers can realistically run on modest consumer hardware---have received disproportionately little attention in the quantization benchmarking literature.
\end{enumerate}

Our project rigorously evaluates \textbf{Qwen2.5-0.5B-Instruct} and \textbf{Llama-3.2-1B-Instruct} (and \textbf{TinyLlama-1.1B} if time permits) in standard FP16, 8-bit, and 4-bit configurations, utilizing official and community releases of \textbf{GPTQ}, \textbf{AWQ}, and \textbf{bitsandbytes NF4/Int8}. We execute six downstream benchmarks (MMLU, ARC-Easy, HellaSwag, WinoGrande, GSM8K, and IFEval) alongside WikiText-2 perplexity, all hosted through one pinned version of EleutherAI's \texttt{lm-evaluation-harness} to guarantee cross-experimental reproducibility.

We address three primary research questions:
\begin{enumerate}[leftmargin=*]
    \item \textbf{Task Sensitivity:} How much does quantization degrade each type of task (factual recall, commonsense reasoning, multi-step arithmetic, and instruction following), across each model family and quantization method?
    \item \textbf{Reliability of the Perplexity Proxy:} Does perplexity actually correlate with task accuracy after quantization, or does it mask severe localized degradation? We quantify this by computing the Spearman rank correlation ($\rho$) between the change in perplexity ($\Delta\text{PPL}$) and the change in accuracy ($\Delta\text{Acc}$) across each task.
    \item \textbf{Practical Bit-Budget Guidelines:} How many bits are sufficient for each task type? We establish a practical threshold calling a setting \textit{acceptable} if accuracy falls by at most 3 percentage points ($\le 3\,\text{pp}$) compared with FP16.
\end{enumerate}

Additionally, we evaluate the hypothesis of whether GSM8K, which relies heavily on exact numerical computation, holds up better under subtle quantization shifts than tasks requiring extensive multi-step commonsense reasoning or reading context.

% ----------------------------------------------------------------------
% 3. BRIEF LITERATURE REVIEW
% ----------------------------------------------------------------------
\section{Brief Literature Review}

\begin{itemize}[leftmargin=*]
    \item \textbf{GPTQ (Frantar et al., ICLR 2023):} Quantizes weights in one shot using approximate second-order inverse Hessian information to correct column errors dynamically, reporting near-lossless 4-bit accuracy on legacy benchmarks.
    \item \textbf{AWQ (Lin et al., MLSys 2024):} Identifies that protecting the top 1\% salient weight channels based on activation magnitude preserves on-device performance without matrix inversion.
    \item \textbf{NF4 \& QLoRA (Dettmers et al., NeurIPS 2023):} Introduces NormalFloat4 (NF4), an information-theoretically optimal quantile quantization for normally distributed weights, integrated into \texttt{bitsandbytes}.
    \item \textbf{Comprehensive Evaluations (Jin et al., ACL Findings 2024):} Evaluated quantized LLMs on ten benchmarks and concluded that perplexity is a reliable proxy for downstream accuracy---the exact conclusion this project re-evaluates under reasoning workloads.
    \item \textbf{Granular Tensor Analysis (Li et al., ICML 2024):} Demonstrated that quantization sensitivity varies dramatically depending on tensor type (weights, activations, KV cache), model scale, and target task.
    \item \textbf{Recent Reasoning Studies (2026):} Oladri et al. discovered that 4-bit models can retain final accuracy while their internal reasoning trajectories shift; Twagirayezu and Mitra proved that INT4 impacts reasoning tasks heterogeneously on 7B models.
\end{itemize}

% ----------------------------------------------------------------------
% 4. TECHNICAL REQUIREMENTS & COMPUTE BUDGET
% ----------------------------------------------------------------------
\section{Technical Requirements and Compute Budget}

\subsection{Hardware Architecture}
\begin{table}[h!]
    \centering
    \renewcommand{\arraystretch}{1.2}
    \begin{tabular}{@{}llp{8cm}@{}}
    \toprule
    \textbf{Resource} & \textbf{Specification} & \textbf{Purpose} \\
    \midrule
    Kaggle GPU & NVIDIA T4 (16 GB VRAM) & Primary execution environment for all GPU evaluation runs (30 GPU-hrs/week quota). \\
    Google Colab & Free Tier T4 GPU & Overflow backup environment if weekly Kaggle quota expires. \\
    Local Laptops & Apple Silicon / CPU & Script authoring, pipeline validation, and testing on small 4-bit checkpoints ($\approx 0.4$ GB). \\
    \bottomrule
    \end{tabular}
\end{table}

\subsection{Software Stack}
\begin{itemize}[leftmargin=*]
    \item \textbf{Core Frameworks:} Python 3, PyTorch, Hugging Face Transformers, Accelerate.
    \item \textbf{Quantized Checkpoint Loaders:} Official and community Hugging Face repositories:
    \begin{itemize}
        \item \texttt{Qwen/Qwen2.5-0.5B-Instruct-GPTQ-Int4} and \texttt{GPTQ-Int8}
        \item \texttt{Llama-3.2-1B-Instruct} community GPTQ-4bit and AWQ-4bit
        \item \texttt{bitsandbytes} NF4 and Int8 runtime quantization
    \end{itemize}
    \item \textbf{Standardized Evaluation:} EleutherAI \texttt{lm-evaluation-harness} (version pinned at 0.4.3).
    \item \textbf{Statistical Analysis:} SciPy (\texttt{scipy.stats.spearmanr}), NumPy, Pandas.
\end{itemize}

\subsection{Compute Feasibility}
The full evaluation grid comprises approximately 40 to 70 evaluation runs. Each run requires between 10 to 40 minutes on an NVIDIA T4 GPU, totaling roughly 15 to 35 GPU-hours. Spreading runs over several weeks fits comfortably within Kaggle's free 30-hour weekly quota.

% ----------------------------------------------------------------------
% 5. EXPECTED OUTCOMES AND WORK PLAN
% ----------------------------------------------------------------------
\section{Expected Outcomes and Work Plan}

\subsection{Deliverables}
\begin{enumerate}[leftmargin=*]
    \item Fully reproducible evaluation pipeline with verified FP16 baselines matching official model cards.
    \item Master evaluation results table across 2--3 models, 3--4 bit-levels, 6 downstream tasks, and WikiText-2 perplexity.
    \item Task-by-task degradation curves and validation of the GSM8K hypothesis.
    \item Spearman rank correlation ($\rho$) matrix determining the fidelity of perplexity as a proxy.
    \item Practical bit-budget guidelines based on the $\le 3\,\text{pp}$ threshold.
    \item Final B.Tech Dissertation and conference-style manuscript.
\end{enumerate}

\subsection{10-Week Gantt Chart (Semester Schedule)}

\begin{center}
\begin{ganttchart}[
    vgrid,
    hgrid,
    x unit=1.1cm,
    y unit chart=0.65cm,
    time slot format=simple,
    bar/.append style={fill=blue!60}
]{1}{10}
\gantttitle{Project Execution Weeks (Oct -- Dec 2026)}{10} \\
\gantttitlelist{1,...,10}{1} \\
\ganttbar{A. Setup \& Dry-runs}{1}{2} \\
\ganttbar{B. FP16 Baselines}{2}{3} \\
\ganttbar{C. Core Runs (MMLU, ARC, etc.)}{3}{4} \\
\ganttbar{Mid-Semester Review}{4}{5} \\
\ganttbar{D. Full Grid (GSM8K, IFEval)}{5}{7} \\
\ganttbar{E. Spearman \& Analysis}{6}{8} \\
\ganttbar{F. Variability (3 Seeds)}{7}{8} \\
\ganttbar{G. Bit-Budget Guidelines}{8}{9} \\
\ganttbar{H. Report \& Paper Draft}{8}{10} \\
\ganttbar{I. Final Viva \& Submission}{9}{10}
\end{ganttchart}
\end{center}

% ----------------------------------------------------------------------
% 6. REFERENCES
% ----------------------------------------------------------------------
\section{References}

\begin{enumerate}[leftmargin=*]
    \item E. Frantar, S. Ashkboos, T. Hoefler, and D. Alistarh, ``GPTQ: Accurate post-training quantization for generative pre-trained transformers,'' in \textit{Proc. International Conference on Learning Representations (ICLR)}, 2023.
    \item J. Lin, J. Tang, H. Tang, et al., ``AWQ: Activation-aware weight quantization for on-device LLM compression and acceleration,'' in \textit{Proc. Machine Learning and Systems (MLSys)}, 2024.
    \item R. Jin, J. Du, W. Huang, et al., ``A comprehensive evaluation of quantization strategies for large language models,'' in \textit{Findings of the Association for Computational Linguistics: ACL}, 2024.
    \item T. Dettmers, A. Pagnoni, A. Holtzman, and L. Zettlemoyer, ``QLoRA: Efficient finetuning of quantized LLMs,'' in \textit{Advances in Neural Information Processing Systems (NeurIPS)}, 2023.
    \item S. Li, X. Ning, L. Wang, et al., ``Evaluating quantized large language models,'' in \textit{Proc. International Conference on Machine Learning (ICML)}, 2024.
\end{enumerate}

\end{document}
